% Français 9118, batch 2 — Gallica views 21–40.
% Pass: Fable 5.1 (claude-fable-5-1), 2026-09-10 — first pass, unchecked against
% the views by a human.
%
% What the twenty views hold. The correspondence continues: four more
% letters of Fourier (one official, three autograph and undated), then the
% three letters of Gauss to Germain that the volume holds, and two leaves
% that are not letters — a copy of Lagrange's note on the 1811 prize memoir,
% and the opening of a « Réponse aux questions de Mlle … » on the lever, in
% a third hand. All in French; the only mathematics set as formulae are
% Lagrange's two equations (v. 39) and the lever proportion (v. 40).
%
%   v. 22    f. 17v–18r Fourier, Académie, 24 juillet 1826: receipt of the
%                       « Remarques sur la nature, les bornes et l'étendue de
%                       la question des surfaces élastiques »; Cauchy named
%                       for a verbal report; the address leaf of Fourier's
%                       15 mars 1824 letter faces it.
%   v. 24    f. 19v–20r Fourier, autograph, « mardi soir »: two tickets, a
%                       box at the Italiens; address leaf of 24 juillet 1826.
%   v. 26    f. 21v–22r Fourier, autograph, « jeudi matin 1er juin »:
%                       Legendre asked him to read a memoir on elastic
%                       surfaces; he will call on Saturday; address leaf of
%                       « mardi soir ».
%   v. 28–29 f. 23v–25r Fourier, autograph, « vendredi matin »: the memoir
%                       read « lundi dernier », commission Laplace, Prony,
%                       Poisson (the séance the 15 mars 1824 letter dates to
%                       9 March); f. 25r a post-scriptum; address leaf of
%                       « jeudi matin » on f. 23v.
%   v. 30–31 f. 25v–26v Fourier, autograph, « vendredi matin »: a cold, and
%                       the election of November (Legendre's, Jussieu's and
%                       Montmorency's votes); address leaf on f. 25v.
%   v. 32    f. 28r     Gauss to « Monsieur Le Blanc », Bronsvic 16 juin
%                       1805: apology for six months' silence, praise of the
%                       new proof for the primes of which 2 is a residue.
%   v. 33    f. 28r     the same leaf, photographed a second time at a larger
%                       scale — not transcribed separately.
%   v. 34    f. 28v–29r the letter's second page: the theorema fundamentale
%                       (art. 131), the residues −1, ±2, the further
%                       demonstrations; Duprat the bookseller's 680 francs;
%                       P.S. about a letter for Orléans.
%   v. 35    f. 29v–30r the address, « Monsieur Le Blanc … chés Mr. Sylvestre
%                       de Sacy … Rue haute feuille N° 11 »; then Gauss,
%                       Bronsvic 20 août 1805: the 1799 memoir sent through
%                       Grégoire, planetary perturbations.
%   v. 37–38 f. 32r–32v Gauss, Gottingue 19 janvier 1808: the move to
%                       Göttingen, the « primices » sent, « Mademoiselle »
%                       and « ma chere amie », the Theoria motus in press.
%   v. 39    f. 34r     « Note communiquée aux commissaires pour le prix de
%                       la surface élastique (décembre 1811) », a copy, with
%                       the note « cette note est de Mr la grange ».
%   v. 40    f. 35r     « Réponse aux questions de Mlle … », a third hand:
%                       the lever with three forces P, S, R; continues on the
%                       verso, in batch 3.
%
% Absent from the output: v. 21 (f. 16v–17r, blank; the 15 mars 1824 letter
% shows through), v. 23 (f. 18v–19r blank), v. 25 (f. 20v–21r blank), v. 27
% (f. 22v–23r blank), v. 36 (f. 30v–31r blank). On v. 31, 32, 37, 38, 39 and
% 40 one of the two facing pages is blank and only the other is transcribed.
% The BnF's stamps and the postmarks are not transcribed.
%
% Folios are the pencilled numbers on the rectos (17–35 across the batch);
% a verso is written as the recto's number with v. No folio was read on
% f. 33 (v. 38, right, blank) beyond the pencil « 33 » itself.
%
% Gauss's French is transcribed as he wrote it: accents are often absent
% (« reponse », « esperance », « cheres ») and are not supplied; his
% signature is « Gauß » with the eszett, set here as « Gauss ». The
% « Monsieur Le Blanc » of 1805 is the name Germain used with Gauss; the
% 1808 letter is addressed to « Mademoiselle ». Lagrange's round d is set
% with \partial; the coefficient of his first equation reads « gEbc » (the
% first letter possibly a 9). Sections are the pass's own, for orientation.
\documentclass[11pt,a4paper]{article}
% The preamble shared by every transcription of the mathematicians' manuscripts in Gallica.
%
% It rests on one idea: **uncertainty compounds, it does not resolve.** A
% transcription that smooths over a doubtful word has destroyed the only thing
% it was made for — a reader must be able to tell, at every point, what was on
% the page and what was supplied. The apparatus macros below are therefore the
% critical apparatus, not decoration.
%
% The same commands are recognised by `scripts/render.mjs`, which produces the
% site's reading view, and by `scripts/tei.mjs`. A macro added here must be
% added there too, or rendering fails loudly — which is the wanted behaviour.
%
% Comments here are English, like the rest of the repository. The documents
% this preamble sets are French, and that is deliberate: see the skills.

\ProvidesPackage{germain}[2026/09/15 Transcriptions of mathematicians' manuscripts in Gallica]

\usepackage{amsmath,amssymb,amsthm}
\usepackage{mathrsfs}
% Kept from the parent preamble so the renderer's diagram support stays
% available; these manuscripts draw no commutative diagrams, and nothing here uses it.
\usepackage{tikz-cd}
\usepackage[english,french]{babel}
\usepackage{fontspec}

% --- Greek ------------------------------------------------------------------
% The text face stays fontspec's default, Latin Modern, which has no Greek:
% XeTeX drops every glyph it lacks with a line in the log and nothing on the
% page, and the Greek words the manuscripts quote vanished from the PDFs. So
% the Greek and Greek Extended blocks get a XeTeX character class of their own,
% and entering or leaving a run of it switches the family to CMU Serif — the
% same Computer Modern design, with polytonic Greek — and back. Only the
% family changes: a Greek word in \textit or \textbf keeps its shape and series.
% The transitions to and from every other class carry the tokens that class
% has towards an ordinary letter, so French spacing before « ; » or inside
% « » still holds after a Greek word. No grouping, so no brace or \endgroup
% can come out unbalanced; maths is left alone, since XeTeX inserts nothing
% there. Kept identical in `exercices.sty`.
\makeatletter
\newfontfamily\germainGreekfont{cmun}[
  Extension=.otf, NFSSFamily=germaingreek,
  UprightFont=*rm, ItalicFont=*ti, SlantedFont=*sl,
  BoldFont=*bx, BoldItalicFont=*bi, BoldSlantedFont=*bl]
\newXeTeXintercharclass\germain@greek
\def\germain@greekrange#1#2{%
  \@tempcnta=#1\relax
  \loop
    \XeTeXcharclass\@tempcnta=\germain@greek
    \ifnum\@tempcnta<#2\advance\@tempcnta\@ne\repeat}
\germain@greekrange{"0370}{"03FF}
\germain@greekrange{"1F00}{"1FFF}
\def\germain@greekprev{\rmdefault}
\def\germain@greekin{%
  \edef\germain@greekprev{\f@family}\fontfamily{germaingreek}\selectfont}
\def\germain@greekout{\fontfamily{\germain@greekprev}\selectfont}
\def\germain@greekjoin{%
  \XeTeXinterchartokenstate=\@ne
  \@tempcnta=\z@
  \loop
    \ifnum\@tempcnta=\germain@greek\else
      \XeTeXinterchartoks\germain@greek\@tempcnta=\expandafter{%
        \expandafter\germain@greekout\the\XeTeXinterchartoks\z@\@tempcnta}%
      \XeTeXinterchartoks\@tempcnta\germain@greek=\expandafter{%
        \the\XeTeXinterchartoks\@tempcnta\z@\germain@greekin}%
    \fi
    \ifnum\@tempcnta<4095\advance\@tempcnta\@ne\repeat}
% babel-french sets its own classes, and saves and restores the interchar
% state, each time a language is selected: join again after it does.
\addto\extrasfrench{\germain@greekjoin}
\addto\extrasenglish{\germain@greekjoin}
\AtBeginDocument{\germain@greekjoin}
\makeatother

\usepackage{geometry}
\geometry{a4paper,margin=25mm,marginparwidth=28mm}
\usepackage{marginnote}
\usepackage{xcolor}
\usepackage[normalem]{ulem}
\setlength{\emergencystretch}{3em}
\usepackage{hyperref}
\hypersetup{colorlinks=true,linkcolor=black,urlcolor=black,pdfborder={0 0 0}}

% --- Batch metadata ---------------------------------------------------------
% Taken as they stand by the rendering script, which shows them at the head of
% the reading view. They fix what the file claims to cover. The macro names are
% the parent project's, so that the scripts stay shared; \folder here means the
% volume's slug — `fr-9115` — and \pages counts Gallica views.

\newcommand{\folder}[1]{\gdef\grothFolder{#1}}
\newcommand{\batch}[1]{\gdef\grothBatch{#1}}
\newcommand{\pages}[2]{\gdef\grothFirst{#1}\gdef\grothLast{#2}}
\newcommand{\foldertitle}[1]{\gdef\grothFoldertitle{#1}}

% The holder's own shelfmark, printed as they write it: « Français 9115 ».
\newcommand{\shelfmark}[1]{\gdef\grothShelfmark{#1}}

% The dating, copied verbatim from the holder's catalogue — never inferred
% here. « XIXe siècle » is the BnF's; a pass that dates a leaf from its content
% says so in a \note{}, as its own inference.
\newcommand{\dating}[1]{\gdef\grothDating{#1}}

% The status notice, printed in the title block and repeated as a diagonal
% watermark on every page of the PDF. The manuscripts are in the public
% domain; what this line declares is not a right but a quality — an
% unverified first pass by a machine. The skills write the line; a file
% without it gets no watermark, so leaving it out is visible in the output.
\newcommand{\watermark}[1]{\gdef\grothWatermark{#1}}

\gdef\grothFolder{?}\gdef\grothBatch{}\gdef\grothFirst{?}\gdef\grothLast{?}%
\gdef\grothFoldertitle{}\gdef\grothShelfmark{}\gdef\grothDating{}\gdef\grothWatermark{}

% --- The résumé -------------------------------------------------------------
\newenvironment{resume}
  {\small\setlength{\parskip}{0.45em}}
  {\par\medskip\hrule\medskip}

% --- Keywords ---------------------------------------------------------------
% In English, closing the modernised reading's résumé: the modern vocabulary
% under which to search for what the leaves construct. The manifest extracts
% them to tag the volume — this line is the single source of the tags.
\newcommand{\keywords}[1]{\par\smallskip\noindent
  {\footnotesize\textsc{Keywords} --- \textit{#1}}\par}

% --- The critical apparatus -------------------------------------------------

% The Gallica view number, in the margin. This is the anchor that lets a
% disagreement over a formula be settled by looking at *one* image rather than
% a volume of 750. The site uses it to turn the facsimile as the transcription
% is scrolled. A view is one scanned image: usually one side of a leaf.
\newcommand{\page}[1]{%
  \marginnote{\footnotesize\color{gray!70}\textsf{v.\,#1}}%
  \label{page:#1}\ignorespaces}

% The BnF's pencilled foliation, where the leaf carries it and a pass can read
% it: « 198r ». This is what the literature cites — « ff. 198r–208v » — and
% the one line that turns a citation into a view. Recorded, never inferred:
% a folio a pass cannot read is not written.
\newcommand{\folio}[1]{%
  \marginnote{\footnotesize\color{gray!50}\textsf{f.\,#1}}\ignorespaces}

% The modernised reading's coarser counterpart to \page{}: which views a
% section is drawn from.
\newcommand{\pagerange}[2]{%
  \marginnote{\footnotesize\color{gray!70}\textsf{v.\,#1--#2}}%
  \ignorespaces}

% Illegible: never guessed.
\newcommand{\ill}{\textcolor{red!60!black}{[\,\ldots\,]}}

% A doubtful reading: the word is given, and flagged as uncertain.
\newcommand{\uncertain}[1]{\uline{#1}}

% An editorial addition — a missing letter, an implied word.
\newcommand{\add}[1]{\textcolor{blue!50!black}{[#1]}}

% Struck out by the writer. What was crossed out is part of the document.
\newcommand{\struck}[1]{\sout{#1}}

% The transcriber's note, on the state of the page or the mathematics.
\newcommand{\note}[1]{\par\smallskip{\small\color{gray!60!black}\textit{[#1]}}\par\smallskip}

% A marginal note of hers, distinct from the body of the text.
\newcommand{\marginal}[1]{\par\smallskip\noindent\hspace{1em}%
  {\small\color{gray!60!black}#1}\par\smallskip}

% --- Title ------------------------------------------------------------------
% The title block is French, like the documents themselves.

\AddToHook{shipout/background}{%
  \ifx\grothWatermark\empty\else
    \begin{tikzpicture}[remember picture, overlay]
      \node[rotate=52, scale=3.4, align=center, text=gray!18,
            font=\sffamily\bfseries]
        at (current page.center) {\grothWatermark};
    \end{tikzpicture}%
  \fi}

\AtBeginDocument{%
  \begin{center}
    {\large\bfseries Manuscrits de mathématiciens (Gallica) ---
      \ifx\grothShelfmark\empty\grothFolder\else\grothShelfmark\fi}\\[2pt]
    {\itshape\grothFoldertitle}\\[4pt]
    {\small vues \grothFirst--\grothLast
      \ifx\grothBatch\empty\else\ (lot \grothBatch)\fi}\\[2pt]
    \ifx\grothDating\empty\else
      {\small\color{gray!60!black}Datation du catalogue : \grothDating}\\[2pt]
    \fi
    \ifx\grothWatermark\empty\else
      {\small\bfseries\color{red!45!gray}\grothWatermark}\\[2pt]
    \fi
    {\footnotesize\color{gray!60!black}Transcription établie d'après les images IIIF de Gallica.
     Source gallica.bnf.fr / Bibliothèque nationale de France.\\
     Les lectures incertaines sont soulignées, les illisibles notées [\,\ldots\,],
     les ajouts entre crochets ; \textsf{v.} numérote les vues de Gallica,
     \textsf{f.} la foliotation au crayon de la BnF.}
  \end{center}
  \vspace{1em}}

\endinput


\folder{fr-9118}
\shelfmark{Français 9118}
\batch{2}
\pages{21}{40}
\dating{1801-1900}
\watermark{Lecture automatique\\non vérifiée}
\foldertitle{Correspondance de Sophie GERMAIN avec les mathématiciens et savants Cauchy, Delambre, Fourier, Gauss, Le Gendre, d'Ansse de Villoison, etc.}

\begin{document}

\section{Lettre de Fourier, Académie royale des sciences, 24 juillet 1826}

\page{22}\folio{17v}
A Mademoiselle

Mademoiselle Sophie Germain

Rue de Braque n°. 4.

à Paris

\note{adresse de la lettre de Fourier du 15 mars 1824, au verso de son second
feuillet, avec deux timbres de la poste}

\folio{18r}
Institut de France

Académie Royale des Sciences.

Paris, le 24 juillet 1826

Le Secrétaire perpétuel de l'Académie.

\note{en-tête imprimé ; la date est à la main}

à Mademoiselle Sophie Germain

Mademoiselle, l'académie a reçu l'ouvrage que vous avez bien voulu lui
adresser et qui est intitulé : \emph{Remarques sur la nature, les bornes et
l'étendue de la question des surfaces élastiques, et équation générale de ces
surfaces}. j'ai l'honneur de vous remercier, au nom de l'académie, de l'envoi
de cet ouvrage. M. Cauchy a été désigné pour en faire un rapport verbal. Le
volume sera déposé dans la bibliothèque de l'Institut

\note{le titre est souligné dans l'original}

j'ai l'honneur Mademoiselle, de vous offrir l'assurance de mon respect.

Bon Fourier.

ci joint un exemplaire de l'analyse des travaux de l'académie année 1825,
partie mathématique.

\note{« Bon » abrège « Baron », en exposant dans l'original ; le post-scriptum
est d'une autre main que la signature}

\section{Billet autographe de Fourier, « mardi soir »}

\page{24}\folio{19v}
A Mademoiselle

Mademoiselle Sophie Germain

\note{adresse de la lettre de Fourier du 24 juillet 1826, au verso de son
second feuillet ; la suite de l'adresse n'est pas écrite}

\folio{20r}
je renouvelle l'expression de mes éternels remercimens pour les témoignages de
bonté et d'amitié que j'ai reçus de mademoiselle germain

j'envoye deux billets dont l'un est converti en billet de centre. jamais on n'a
montré autant d'empressement et il y a plus d'un mois que monsieur cuvier et
moi avons reçu des demandes sans nombre. mais n'eussé-je qu'un seul billet
\ill{} certainement pour mademoiselle germain.

\note{entre « billet » et « certainement » deux mots d'une écriture très
rapide n'ont pu être lus}

je suis encor bien incertain de savoir si je pourrai profiter de l'offre très
obligeante et très aimable concernant la place de la loge. car le matin nous
aurons une longue séance. mais si je vais le soir aux italiens ce seroit
seulement pour occuper une place dans la loge et je ne pourrois y aller qu'à
huit heures et demie. je prie bien instamment mademoiselle germain de disposer
de cette loge et je renvoye le billet. mais si elle a la bonté d'insister à
cet égard je désire qu'elle veuille bien se contenter de me renvoyer par le
porteur un bon pour une place. j'aurai \uncertain{été} assez heureux pour en
profiter. mille respects et \uncertain{vrais} remercimens

Jh. Fourier.

mardi soir.

\note{« Jh » abrège « Joseph », en exposant dans l'original}

\section{Lettre autographe de Fourier, « jeudi matin 1\textsuperscript{er} juin »}

\page{26}\folio{21v}
Mademoiselle

Mademoiselle Sophie germain

rue de brac n° 4.

au marais

Paris.

\note{adresse du billet « mardi soir », au verso de son second feuillet}

\folio{22r}
Mademoiselle,

Monsieur le gendre a bien voulu m'engager de votre part à prendre connoissance
d'un mémoire sur les propriétés des surfaces élastiques. j'ai lu fort
attentivement cet écrit, et j'y ai trouvé de nouvelles preuves de l'importance
et du succès de vos recherches sur cette question difficile. je me propose
d'avoir l'honneur de me rendre chez vous après demain samedi à huit heures et
demie du soir, et de vous rendre compte de mes réflexions sur l'objet de ce
mémoire. cette heure m'a été indiquée comme vous étant la plus commode. si
vous préfériez une autre heure, ou un autre jour, je vous prie d'avoir la bonté
de le faire dire au porteur de cette lettre. je serai très empressé de m'y
conformer, et dans le cas où l'on ne me donneroit de votre part aucune
indication différente, j'aurai l'honneur de me présenter samedi.

\note{« le gendre » est écrit en deux mots : Legendre}

agréez mademoiselle l'hommage du respect de votre très humble et obéissant
serviteur

Fourier.

jeudi matin 1\textsuperscript{er} juin.

\section{Lettre autographe de Fourier, « vendredi matin », avec post-scriptum}

\page{28}\folio{23v}
Mademoiselle

Mademoiselle Sophie germain

rue de Braque

maison de monsieur geoffroy médecin

au marais

Paris.

\note{adresse de la lettre du « jeudi matin 1\textsuperscript{er} juin », au
verso de son second feuillet}

\folio{24r}
mademoiselle

Je regrette extrêmement de n'avoir pu répondre aussi promptement que je
l'aurois désiré au sujet du mémoire de mathématique que vous nous avez envoyé.
je me suis acquitté fidèlement de la commission que vous m'aviez donnée en
m'adressant cette pièce. Mr cuvier étoit chargé lundi dernier de la lecture de
la correspondance. je l'ai prié de présenter votre mémoire et j'en ai indiqué
l'objet. après la lecture on a nommé m.m. laplace prony et poisson
commissaires. j'insisterai autant qu'il sera nécessaire pour qu'ils fassent
le rapport que vous désirez. si m. poisson a le dessein de montrer
quelqu'opposition au résultat de vos recherches il ne pourra s'empêcher de
céder à l'autorité de l'expérience que personne ne sait mieux consulter que
vous. autant que j'ai pu prendre connoissance de la discussion dont vous vous
êtes occupée il m'a paru que vous mettez dans tout son jour l'insuffisance de
l'hypothèse théorique dont on a voulu déduire l'équation du
4\textsuperscript{eme} ordre que vous avez trouvée.

\note{la lecture « lundi dernier » et la commission Laplace, Prony, Poisson
sont celles que la lettre officielle du 15 mars 1824 (vue 20) rapporte à la
séance du 9 mars ; ce billet du « vendredi matin » en seroit le compagnon
autographe, du 12 mars 1824 — inférence de la présente lecture, non écrite sur
le feuillet}

\page{29}\folio{24v}
je n'aurois pu concourir moi même à l'examen et au rapport du mémoire sans me
détourner des occupations instantes dont je me trouve chargé. toutes les
personnes présentes à la séance ont entendu avec le plus grand intérêt
l'annonce de votre mémoire. la difficulté du sujet, la célébrité des auteurs
qui l'ont traité et votre nom ne pouvoient manquer d'exciter l'attention. nous
nous en sommes entretenus avec plusieurs personnes à l'académie et chez mr de
laplace.

je vous remercie mademoiselle des nouvelles marques d'intérêt que vous me
donnez en vous occupant de ma santé et de mes travaux. c'est une obligation
\uncertain{nouvelle} que celle des \ill{} vous publiez et les personnes dont
j'estime le plus les suffrages sont ceux que je crains le plus d'avoir pour
\struck{\ill{}} \uncertain{censeurs}.

\note{un mot après « celle des » et le mot biffé avant « censeurs » n'ont pu
être lus ; « censeurs » est écrit par-dessus la biffure}

j'aurois préféré de vous rendre compte de vive voix au sujet de la
présentation de votre mémoire et je profiterai d'une autre occasion pour vous
en parler je suis présentement \uncertain{retenu} par des occupations beaucoup
moins agréables

agréez mademoiselle avec l'hommage de mes remercimens celui de mon respect.

Jh. fourier.

vendredi matin.

\folio{25r}
P. S.

le procès verbal que j'ai rédigé contient la mention de la lecture de votre
mémoire la lettre par laquelle je vous informe du nom des commissaires ne vous
est \ill{} parce qu'on a coutume de les expédier \ill{} que quand le procès
verbal a été lu et adopté.

Fr.

\note{deux taches d'encre couvrent un mot après « mémoire » et un mot après
« expédier » ; le mot après « ne vous est » n'a pu être lu}

\section{Lettre autographe de Fourier, « vendredi matin », sur une élection}

\page{30}\folio{25v}
Mademoiselle

Mademoiselle Sophie germain

rue de braque n° 4. au marais

Paris.

\note{adresse de la lettre du « vendredi matin » qui précède, au verso du
feuillet portant son post-scriptum}

\folio{26r}
mademoiselle,

je ne puis aller vous exprimer combien je suis reconnoissant de l'intérêt que
vous m'accordez et de la grace parfaite avec laquelle vous l'exprimez. les
personnes que vous aimez et que vous protégez ne doivent pas être
malheureuses. je me permettrai de vous recommander de ne point sortir : car
l'air est très froid, et un grand nombre de personnes sont fort incommodées.
j'étois revenu à pied mardi soir du fauxbourg St honoré j'ai été saisi d'un
rhume violent qui m'a causé des douleurs vives dans tout le corps la médecine
\uncertain{adoptant} le langage de la géométrie appelle cette indisposition
une courbature. la mienne étoit certainement d'un degré très élevé. enfin elle
a cessé entièrement et j'ai pu sortir.

je ne puis douter maintenant que le vœu du plus grand nombre de mes collègues
soit de me choisir, et celui de mes concurrens qui se flatte le plus est dans
une grande erreur mais il a recours à tant d'artifices qu'il y auroit de

\page{31}\folio{26v}
l'imprudence à ne pas les redouter.

mr desfontaines m'a dit que mr legendre s'étoit entretenu avec lui de cette
élection et que sans disconvenir de l'intérêt qu'il prenoit à mr D. il l'avoit
assuré que dans tous les cas possibles il me donneroit son suffrage mr
desfontaines en paroît convaincu. je ne doute pas mademoiselle que votre
démarche et permettez moi de le dire votre éloquence ne l'ayent touché. un
suffrage que je vous devrois a encore plus de prix à mes yeux. celui de mr de
jussieu est très honorable par lui même et je ne doute pas que vous ne l'ayez
déterminé. l'élection n'aura lieu qu'au commencement de novembre, je
\uncertain{suis sûr} qu'à cette époque mr de jussieu \uncertain{sera} encor à
la campagne. j'apprends avec peine l'absence de mr de montmorency car son avis
m'auroit été favorable. enfin les dieux en décideront. mais ce qui est
indépendant des dieux ce sont mes sentimens de reconnoissance et de respect
je vous prie d'en agréer l'hommage.

\note{« mr D. » : l'initiale seule est écrite}

vendredi matin. fourier.

\section{Lettre de Gauss à « Monsieur Le Blanc », Brunswick, 16 juin 1805}

\page{32}\folio{28r}
Bronsvic 16 Juin 1805.

Monsieur

Il me faut Vous demander mille fois pardon d'avoir laissé six mois sans
réponse l'obligeante lettre, dont Vous m'avés honoré. Certainement je me
serais empressé, de Vous témoigner tout de suite, combien m'est cher l'intérêt
que Vous prenés aux recherches, auxquelles j'ai devoué la plus belle partie de
ma jeunesse, qui ont été la source de mes jouissances les plus délicieuses et
qui me seront toujours plus cheres qu'aucune autre science. Mais je me
flattais de tems en tems, de pouvoir gagner assés de loisir, pour mettre en
ordre et Vous communiquer par écrit une ou l'autre de mes autres recherches
arithmetiques pour Vous rendre en quelque sorte le plaisir que Vous m'avés
fait par Vos communications. Mon esperance a été vaine : ce sont surtout mes
occupations astronomiques, qui à present absorbent presque tout mon tems. Je
me reserve pourtant de m'entretenir avec Vous des mysteres de mon Arithmetique
cherie, aussitôt que je serai assés heureux d'y pouvoir retourner.

J'ai lu avec grand plaisir les choses que Vous m'avés bien voulu communiquer ;
je me felicite, que l'Arithmetique acquiert en Vous un ami aussi habile.
Surtout Votre nouvelle demonstration pour les nombres premiers, dont 2 est
residu ou non residu m'a extremement plu ; elle est très fine, quoique elle
semble être isolée et ne pouvoir s'appliquer à d'autres nombres.

\note{la vue 33 reproduit le même feuillet 28r une seconde fois, à plus
grande échelle ; elle n'est pas transcrite à part. « Bronsvic » est la
graphie de Gauss pour Brunswick ; l'accent manque souvent sur « e » dans sa
main (« reponse », « esperance », « arithmetiques ») et n'est pas suppléé ;
« cheres », « très fine » comme écrit. La lettre est adressée à Monsieur Le
Blanc, nom sous lequel Germain écrivoit à Gauss (voir l'adresse, f. 29v)}

\page{34}\folio{28v}
J'ai très souvent consideré avec admiration l'enchainement singulier des
verités arithmetiques. Par exemples le theoreme que je nomme fondamental (art.
131) et les theoremes particuliers concernans les residus $-1$, $\pm 2$,
s'entrelacent à une foule d'autres verités, où l'on ne les auroit jamais
cherchés. Outre les deux demonstrations que j'ai donnés dans mon ouvrage je
suis en possession de deux \struck{ou} trois autres, qui du moins ne le cedent
pas à celles-la en question d'elegance.

Je remarque avec beaucoup de regret, que les autres occupations ou je suis
engagé ne me permettent point du tout de me livrer à present à mon
\uncertain{amour} pour l'Arithmetique. ce ne sera peut etre qu'après
plusieurs années que je pourrai penser à la publication de la suite de mes
recherches, qui rempliront aisement un ou deux volumes \struck{de la même}
semblables au premier. Mais je croirois n'avoir pas assés vecu, si je mourrais
sans avoir achevé toutes les recherches intéressantes, \struck{\ill{}}
auxquelles je me suis une fois livré. Au reste, chés nous en Allemagne, la
publication d'un tel ouvrage a ses difficultés ; quoiqu'on en dise le gout
pour les Mathematiques pures, si l'on cherche de la profondeur, n'est pas trop
general, nos libraires ne se melent gueres de ces sortes de livres, et je ne
suis pas assés riche pour faire à mes frais l'impression et me soumettre à la
malhonnetteté des libraires etrangers, comme il m'est arrivé à l'occasion du
premier volume. Un M. Duprat, par exemple, qui se nomme libraire pour le bureau
de longitude à Paris, a reçu de moi il y a presque 3 ans, des exemplaires pour
la valeur de 680 francs, mais jamais je n'ai reçu un sous de lui, et il ne
s'est pas même donné la peine de repondre à mes lettres. Peut-être Vous
pourriés me donner des renseignemens, par quels moiens on pour-

\note{« ou » (« ou je suis engagé ») sans accent, comme écrit ; le mot après
« intéressantes » est biffé et illisible ; « ne s'est » est ajouté dans
l'interligne au-dessus d'un mot biffé ; « théorème fondamental (art. 131) »
renvoie à l'article 131 des \emph{Disquisitiones arithmeticae}, la loi de
réciprocité quadratique ; « le premier volume » est ce même livre}

\folio{29r}
rait engager cet homme à faire son devoir.

Agréés, Monsieur l'expression, de ma haute consideration

Ch. Fr. Gauss.

P. S. Oserai-je Vous prier de dépêcher à Orleans la lettre ci-jointe ?

\note{la signature est écrite « Gauß », avec l'eszett ; le post-scriptum est
au bas du feuillet, sous un espace blanc}

\page{35}\folio{29v}
A Monsieur

Monsieur Le Blanc

à Paris

à remettre chés Mr. Sylvestre de Sacy membre de l'institut national. Rue
haute feuille. N° 11

\note{adresse de la lettre du 16 juin 1805, au verso de son second feuillet,
écrite en travers ; la mention « à remettre chés Mr. Sylvestre de Sacy » est
en petits caractères, le long du pli ; cachet de cire}

\section{Lettre de Gauss, Brunswick, 20 août 1805}

\folio{30r}
Bronsvic 20 Aout 1805.

Je profite de la complaisance de Mr. Gregoire, pour Vous offrir, avec beaucoup
de remercimens pour toutes les communications de Votre derniere lettre, un
exemplaire d'un petit memoire, que j'ai publié en 1799, et qui probablement
Vous sera encor inconnu. Vous souhaitiés de savoir tout ce que j'ai écrit en
latin ; cette piece est la seule outre mes recherches arithm : , et en même
tems \struck{la premiere} celle qui a paru la premiere, quoique alors
l'impression de mes Disqu : eut été portée au-delà de la moitié.

Je suis à present occupé à perfectionner quelques methodes \struck{ayant}
nouvelles rapport aux calculs des perturbations planetaires : celles-ci et les
methodes dont je me suis servi pour calculer les elemens elliptiques des
differentes nouvelles planetes, fourniront probablement les materiaux pour
\struck{mon} premier ouvrage.

\note{« nouvelles » est écrit dans l'interligne au-dessus de « ayant » biffé,
mais la place du mot dans la phrase reste celle de « ayant » ;
« elliptiques » est ajouté dans l'interligne au-dessus de « des » ; le
mémoire de 1799 est la thèse de Helmstedt sur le théorème fondamental de
l'algèbre — identification de la présente lecture, non écrite sur le
feuillet}

Je Vous salue cordialement

Ch. Fr. Gauss

\section{Lettre de Gauss, Göttingen, 19 janvier 1808}

\page{37}\folio{32r}
Gottingue ce 19 Janvier 1808

En Vous remerciant de tout mon Cœur pour Votre derniere lettre et les
interessantes communications que Vous m'y faites, Mademoiselle, je Vous prie
mille fois pardon, d'y repondre aussi tard. Cette negligence est pour la plus
grande partie \struck{la} une suite des changemens qui se sont \struck{passés}
faits dans ma situation. J'ai changé ma demeure, pour accepter la place de
professeur d'astronomie a Gottingue qu'on m'avait offerte depuis longtems. Je
ne Vous dirs rièn des circonstances facheuses qui m'ont enfin determiné à
faire ce pas ni des nouvelles tracasseries auxquelles je me trouve exposé ici :
j'espere que l'interposition de l'institut ou j'ai eu recours y mettra fin.
Ne contemplons à présent que la belle perspective que j'ai, de pouvoir avec
plus d'aisance du moins dans la suite, \struck{de} veiller à mes travaux sur
tout arithmetiques et de les publier successivement dans les memoires de la
Société de Gottingue. J'ai le plaisir de Vous en envoier \struck{\ill{}} les
primices lesquelles comme j'espere Vous feront quelques plaisir. Vous me
pardonnerez que cette fois je ne puis m'étendre davantage sur les belles
demonstrations de mes theoremes arithmetiques. J'admire la sagacité avec
laquelle Vous avez pu en si peu de tems y parvenir. J'espere de pouvoir
bientot publier toute la theorie, dont ces propositions elegantes font partie,
avec une foule d'autres choses. Que mes occupations arithmetiques me rendent
heureux dans un tems ou je ne Vois autour de moi que le malheur et le
desespoir ! Ce ne sont que les sciences, le sein de sa famille et la
correspondance avec ses amis cheris ou l'on puisse se dedommager et se
reposer de l'affliction generale.

\note{« faits » est écrit dans l'interligne au-dessus de « passés » biffé ;
« dirs » et « rièn » comme écrit ; deux mots biffés avant « les primices »
n'ont pu être lus ; « sur tout » en deux mots}

L'ouvrage sur le calcul des orbites des planetes, dont je Vous ai parlé dans
ma dernière lettre, est enfin sous presse. J'espere qu'il sera achevé dans
quelques mois. Je n'ai pas redouté la peine de le traduire en \emph{latin},
afin

\note{« latin » est souligné dans l'original}

\page{38}\folio{32v}
qu'il puisse trouver un plus grand nombre des lecteurs.

Soies toujours aussi heureuse, ma chere amie, que Vos rares qualités d'esprit
et de cœur le meritent, et continués de tems en tems de me renouveller la
douce assurance, que je puis me compter parmi le nombre de Vos amis : titre
duquel je serai toujours orgueilleux.

Ch. Fr. Gauss.

\section{Note de Lagrange communiquée aux commissaires du prix, décembre 1811}

\page{39}\folio{34r}
Note communiquée aux commissaires pour le prix de la surface élastique
(décembre 1811.)

L'équation fondamentale pour le mouvement de la surface vibrante ne me paraît
pas exacte, et la manière dont on cherche à la déduire de celle d'une lame
élastique en passant d'une ligne à une surface, ne paraît peu juste. lorsque
les $z$ sont très petits l'équation se réduit à
\[
\frac{\partial^2 z}{\partial t^2}
+ gEbc\left(\frac{\partial^6 z}{\partial x^4 \partial y^2}
+ \frac{\partial^6 z}{\partial y^4 \partial x^2}\right) ;
\]
mais en adoptant comme l'auteur $\frac{1}{r} + \frac{1}{r'}$ pour la mesure
de la courbure de la surface, que l'élasticité tend à diminuer, et à laquelle
on la suppose proportionnelle, je trouve dans le cas des $z$ très petits une
équation de la forme
\[
\frac{\partial^2 z}{\partial t^2}
+ k^2\left(\frac{\partial^4 z}{\partial x^4}
+ 2\frac{\partial^4 z}{\partial x^2 \partial y^2}
+ \frac{\partial^4 z}{\partial y^4}\right) = 0
\]
qui est bien différente de la précédente

(cette note est de Mr la grange)

\note{copie d'une main qui n'est pas celle de Germain ; « ne paraît peu
juste » comme écrit ; le coefficient de la première équation se lit « gEbc »,
la première lettre pouvant aussi être un 9 ; la mention « cette note est de
Mr la grange » est entre crochets dans l'original, d'une écriture plus
rapide ; l'équation « fondamentale » critiquée est celle du mémoire de 1811}

\section{« Réponse aux questions de Mlle », mécanique du levier}

\page{40}\folio{35r}
Réponse aux questions de Mlle \ill{}

\note{le nom qui suit « Mlle » est réduit à un trait ondulé ; le titre est
souligné ; la main n'est pas celle de Germain ; un croquis dans la marge
gauche montre une ligne horizontale $A$ $C$ $B$, avec en $A$ une flèche $P$
dirigée vers le bas, en $C$ une flèche $S$ dirigée vers le haut, et en $B$ un
trait descendant vers $R$}

Pour faire sentir et evanouir le Paradoxe apparent qui resulte de la
proposition de M. Monge, nous allons resoudre plus generalement le probleme
qu'il paroît proposer

\note{« proposition » est abrégé « propoon » dans l'original}

aux 3 points $A$. $C$. $B$ d'une ligne $AB$ sont appliquées les 3 puissances
$P$ poussant vers $A$. $S$ vers $C$ et $R$ vers $B$, et elles sont en
equilibre, on demande les rapports de ces 3 puissances.

Les principes ordinaires de la mecanique suffisent pour voir que puisque ces
3 puissances sont en equilibre, si nous prenons $B$ pour le point d'appuy la
puissance $S$ sera a $P$ dans le rapport inverse de leurs distances au point
d'appuy $B$. c'est a dire que $S : P :: AB : CB$, ou $S \times CB = P \times
AB$.

\note{« principes » est abrégé « ppes » dans l'original}

Supposons que $S$ reste dans sa position, mais que la puissance $P$ s'eloigne
de $B$, il n'y aura plus equilibre a moins que $P$ ne diminue en \struck{p}
même proportion que la distance $AB$ augmentera, afin que le produit de $P
\times AB$ egale toujours le produit constant

\note{le texte se poursuit au verso, dans le lot suivant}

\end{document}
